\begin{table}[htbp]
\centering
\caption{Entry Density: Border vs.\ Interior, by Licensing Regime}
\label{tab:entry_density_mechanism}
\small
\begin{tabular}{p{6.4cm}ccc}
\toprule
& Coefficient & $p$-value & $n$ (counties) \\
\midrule
Colorado + Washington (capped), dispensaries per 1,000, full controls & $-0.156$ & $0.043$ & 103 \\
Montana (uncapped), dispensaries per 1,000, full controls & $+0.142$ & $0.106$ & 23 \\
\quad Montana, no controls & $+0.178$ & $0.035$ & 23 \\
\bottomrule
\end{tabular}
\begin{minipage}{0.95\textwidth}
\vspace{0.3em}
\footnotesize
\textit{Notes:} Outcome is dispensary count per 1,000 residents; ``full controls'' = log(population)
(+ state fixed effects for the CO+WA specification, which pools both states). Border counties in the two
capped states have significantly \emph{fewer} dispensaries relative to population than interior counties,
not merely a null -- the demand implied by the revenue results is not being met by additional entry.
Montana's much smaller sample (23 counties, 8 border) limits precision, but the coefficient is
directionally consistent with the opposite prediction: entry responds to demand where it is not capped.
\end{minipage}
\end{table}
